\clearpage\thispagestyle{empty}
\pdfbookmark[0]{Quellen zu Mondbeobachtungen (April 1822)}{moonapril1822-sources}
\begin{center}\textbf{Quellen zu Mondbeobachtungen (April 1822)}\end{center}
\noindent Diese KI-Übertragung ins heutige Deutsch folgt der zweiten Mitteilung mit der Überschrift \textit{Mondsbeobachtungen} in \href{https://gdz.sub.uni-goettingen.de/id/PPN236007467}{Gauß’ \textit{Werken}, Bd.~VI} (Göttingen, 1874), S.~441. H. C. Schumacher veröffentlichte sie zuerst unter derselben Überschrift in den \textit{Astronomischen Nachrichten}, Bd.~I, \href{https://api.digitale-sammlungen.de/iiif/image/v2/bsb10049426_00103/full/1800,/0/default.jpg}{Nr.~12}, April 1822, Sp.~\href{https://api.digitale-sammlungen.de/iiif/image/v2/bsb10049426_00109/full/1800,/0/default.jpg}{189--190}, und zeichnete mit S. Das \href{https://api.digitale-sammlungen.de/iiif/image/v2/bsb10049426_00005/full/1800,/0/default.jpg}{Titelblatt} des gebundenen Bandes aus der Bayerischen Staatsbibliothek nennt Altona 1823 und die Hammerich- und Heineking’sche Buchdruckerei. Schumacher teilt die Göttinger Messungen von Gauß mit und ergänzt Beobachtungen von Kapitän v. Caroc in Kopenhagen.
\par\smallskip\noindent Für 309 Mayer am 3. März steht in der Werkausgabe $-13\quad3.60$, im Erstdruck $-2\quad46{,}94$. Für 371 Mayer am 4. März steht in der Werkausgabe $-2\quad24\;57$ ohne Dezimalzeichen, im Erstdruck $-2\quad24{,}57$. In der Werkausgabe bleiben die Differenzen in den vier Mondzeilen leer, wo im Erstdruck Striche stehen. Sie verwendet hochgestellte m und s, wo der Erstdruck einfache und doppelte Striche verwendet. KI-Übertragung und KI-Satz.
\clearpage
